\chapter{Introduction and background}
\label{ch:intro}

\begin{instructions}
Restructured 5 Nov. This chapter now holds the introduction and the fungal background: the problem
(1.1), previous research (1.2), what is known about electrical signals in fungi (1.3, from the old Chapter 2), the missing
signal vocabulary (1.4), the summer project at UQ as separate prior work (1.5), the gaps (1.6), the research questions as
they stand (1.7), contributions (1.8), scope (1.9) and the overview (1.10). Definition is 10\,\% of the mark and
Background 20\,\%, shared with Chapter~\ref{ch:litreview}. About 8--9 pages.
\end{instructions}

\begin{figure}[htbp]
\centering
\figplaceholder{5cm}{Figure 1.1 (T1, bundle): one M2\_aug channel at three zooms --- the whole 721\,h, one day, and
30 minutes --- with the zoomed region boxed at each level. The ``1,000-hour problem'' in one picture.}
\caption{One electrode of the primary oyster recording at three zooms.}
\label{fig:threezooms}
\end{figure}

\section{Problem space}

Bio-electrical potential has been observed and measured in mycelium networks for the past five decades. Early work by Slayman, Long and Gradmann \cite{slayman} measured the potential differences exhibited by the fungus \emph{Neurospora crassa} and found that the signal demonstrates evidence of ``action potential''-like spikes. The measured action potentials in the hyphae of the mycelium exhibited characteristics not unlike those measured in cardiac fibres, although at four orders of magnitude larger, with rapid polarisation and slow depolarisation behaviour. Further work on measuring these signals \cite{adamatzky2023} \cite{dehshibi} \cite{mayne} has found multiple families and scales of oscillatory behaviour across different species of fungi, with unique but distinctive patterns of activity.

Bio-electric signalling behaviour across many domains is increasingly being recognised as a fundamental basis for development, regeneration and communication within biological networks \cite{zhang}. Understanding the mechanisms and purposes of signalling behaviour can lead to significant discoveries in creating systems to control the behaviour of cellular collectives and reduce undesirable effects. Bioelectric stimulation has been used to trigger transient permeabilization to enhance local uptakes of drugs that kill dividing cancer cells, to generate reversible nano-pores in membranes to treat skin lesions, and to modify membrane potentials of brain tissue to reverse abnormal brain and eye morphologies created by ethanol exposure \cite{nuccitelli}.

There are many open questions that the existence of these bio-electrical signals in mycelium networks propose. More notably, there arise the questions of determining the meaning of the signals in the context of the growth and functions of a network, using this to predict responses of a network based on current recordings or stimuli, and leveraging this knowledge to manipulate mycelium growth and behaviour. There is also research into the prospect of using mycelium networks as a medium for information exchange via signal propagation \cite{mayne} \cite{phillips}.

Parallel to the field of bio-electrics is time-series analysis, which is a collection of techniques and algorithms for extracting information from measured data in time. Time-series analysis is a significant branch of both mathematics and statistics, useful across the domains of finance, business, energy and biological systems, to name a few. As scientific research collects more data from increasingly complex dynamical systems, the methods have adapted to describe them. Traditional methods of decomposition and smoothing algorithms have moved towards more flexible methods that involve machine learning and feedback-based approaches, leveraging the increase in processing power and memory available today \cite{cabauatan} \cite{gao} \cite{milvus}.

The models for time-series analysis are constantly improving and changing, and their applications to different fields are always lagging slightly behind. The application of some of the more modern techniques of time-series analysis to the electrical potentials recorded for mycelium networks is severely under-examined and could present novel ideas about their behaviour and functions. Also missing is a standard theoretical framework for the identification and classification of the families of subsequences seen in mycelium signals across species.

\begin{instructions}
Add one paragraph in your own words on the motivation as you experienced it, reworked from your note to the demo
session (20 Oct). Paraphrase of that note: you were handed about 1,000 hours of raw recordings with no field standard
for what to look for; exploratory algorithms gave wildly different answers that could not be compared; each new
combination of ideas was a one-off script that was then lost; fungi have no defined alphabet, let alone a language.

Back it with numbers:
\begin{itemize}
  \item 721\,h + 282.5\,h $\approx$ 1,003\,h across the two primary oyster recordings, M2\_aug and M2\_concat, each
        16 channels;
  \item OASIS marks 61--65\,\% of M2\_aug's time as events, while drop\_detection\_v1 finds about a thousand;
  \item a seed bundle whose generator scripts were deleted, so it cannot be regenerated
        (\src{DATA/library_seed/PROVENANCE.md}).
\end{itemize}
Point to Figure~\ref{fig:threezooms}.
\end{instructions}

\section{Scope of previous research}

Research into the bio-electric signals of mycelium networks has developed across a small number of broad areas. The earliest and most established of these is the recording and characterisation of the signals themselves, with studies across multiple species documenting the waveform morphology, amplitude, and oscillatory periods of the measured potentials \cite{slayman,adamatzky2023,dehshibi}. A parallel body of work has focused on the analysis of these recordings, using spike-detection algorithms and simple frequency methods to extract and describe the behaviour of individual events within the signal \cite{adamatzky2023,dehshibi,buffi}.

More recently, attention has turned to questions of what these signals might mean and what they might be used for. Some work has explored the propagation of signals across the mycelium substrate and the feasibility of using mycelium networks as a medium for data transmission, finding that the substrate exhibits low-pass filter characteristics relevant to this application \cite{mayne,phillips}. Others have taken a more interpretive approach, treating the spike-trains as a form of language and performing linguistic analysis to compare the complexity of signals across species \cite{adamatzky2022}. There is also a smaller body of work examining the biological response of the signals to external stimuli such as injected current and chemical agents, with the aim of understanding whether and how the signals can be manipulated \cite{buffi,olsson}.

Sections~\ref{sec:litfungi} and~\ref{sec:litvocab} review the first of these areas in more detail. The time-series methods are reviewed in Chapter~\ref{ch:litreview}.

\section{Analysis of electrical signals in fungi}
\label{sec:litfungi}

Only recently has the idea of communication within fungi networks via neuron-like transmission become popularised \cite{hunter} \cite{buffi2025review}. The earliest recorded study on the electrical signals in fungi was recorded by Clifford Slayman, Scott Long and Dietrich Gradman in 1976 \cite{slayman}, after which there was a significant stall in the progress made in this field \cite{buffi2025review}. From inserting 1--3 micro-electrodes directly into the hyphae (mycelium cells) of a small inoculum of \emph{Neurospora crassa} on the surface of scratched cellophane, they found quasi-sinusoidal fluctuation of 10--20\,mV amplitude, persisting for varying intervals of time, `\emph{from 10-15 minutes to 1 hour or more, and which has a constant period of 34 minutes}'. They describe the oscillatory behaviour as a sharp depolarisation phase followed by a slow re-polarisation phase, a trend also observed in similar studies on electrical signals \cite{adamatzky2023} \cite{springer}. Multiple scales of these behaviours have been observed in mycelium, and these behaviours are often species-specific.

Efforts have been made to begin to understand the functions of these signalling characteristics. More recent work has examined additional species and recording configurations in order to characterise these signals more fully and test their response to external stimuli. Adamatzky et al. \cite{adamatzky2023} recorded multi-scale electrical spiking in \emph{Schizophyllum commune} using a multi-electrode array, modelling the signal using the FitzHugh--Nagumo system and detecting and characterising individual spike events via a simple thresholding algorithm. Dehshibi and Adamatzky \cite{dehshibi} extended this spike detection approach, implementing a more complex Morse wavelet transform combined with a signal envelope threshold. They evaluated their results with entropy-based complexity analysis, finding that spike family characteristics are different amongst scales, and comparable to the complexity of modern languages, similar to the findings of another study by Adamatzky \cite{adamatzky2022}. Buffi et al. \cite{buffi} grew mycelium across two electrodes inside a Faraday cage and used the Short-Term Fourier Transform to extract spike frequency content, showing that the fungicide cycloheximide had a significant dampening effect on frequency activity. A study by Olsson and Hansson \cite{olsson} found that action-potential frequency in \emph{Pleurotus ostreatus} and \emph{Armillaria bulbosa} increases in response to injected negative current, indicating a difference in the behaviour of these action potentials from that observed in animal nerve cells.

There is also research into the propagation of these signals across the mycelium substrate. Mayne et al. \cite{mayne} measured the propagation of electrical signals between electrodes and found that the mycelium substrate exhibits low-pass filter characteristics, relevant to its proposed use as a medium for data transmission \cite{phillips}.

The recording and analysis parameters across these key studies are summarised in Table~\ref{tab:params}. It is worth noting that sampling rates are not consistently reported across the literature, and where they are, resampling to 1\,Hz appears to have become an implicit standard, as used notably in Adamatzky's language analysis \cite{adamatzky2022}. It can be seen that a variety of fungi are studied with a variety of techniques, suggesting that there is still a large degree of exploration in this field. The goals of each paper vary significantly, with more emphasis often on methods for producing and detecting significant events rather than classifying or explaining their functionality.

\begin{table}[htbp]
\caption{Comparison of recording and analysis parameters across selected studies. Sampling rates are taken from full paper methods sections where stated.}
\label{tab:params}
\footnotesize
\begin{tabularx}{\textwidth}{p{2.1cm} p{2.3cm} p{2.2cm} p{1.7cm} p{2.4cm} X}
\toprule
\textbf{Study} & \textbf{Species} & \textbf{Electrode type} & \textbf{Sample rate} & \textbf{Preprocessing} & \textbf{Analysis / key finding}\\
\midrule
Slayman et al.\ \cite{slayman}, 1976 & \emph{Neurospora crassa} & 1--3 intracellular microelectrodes & Not stated & Electrometer amplification & Oscillogram inspection; 34-min oscillation, action-potential spikes\\
\addlinespace
Adamatzky et al.\ \cite{adamatzky2023}, 2023 & \emph{Schizophyllum commune} & Multi-electrode array & 1\,Hz (600 per s, averaged) & Resampled to 1\,Hz & FitzHugh--Nagumo modelling; species-specific spike families\\
\addlinespace
Dehshibi \& Adamatzky \cite{dehshibi}, 2021 & \emph{Pleurotus ostreatus} & Iridium-coated stainless-steel subdermal & 1\,Hz (600 per s, averaged) & Morse wavelet + envelope threshold & Morse wavelet spike detection; entropy complexity analysis\\
\addlinespace
Buffi et al.\ \cite{buffi}, 2025 & \emph{Fusarium oxysporum} & 2-electrode (2\,mm), Faraday cage & 17\,Hz & NaN values removed, low-pass to 4\,Hz, subsampled to 8\,Hz & STFT spike frequency extraction; cycloheximide dampens activity\\
\addlinespace
Adamatzky \cite{adamatzky2022}, 2022 & \emph{Omphalotus nidiformis}; \emph{Flammulina velutipes}; \emph{S.\ commune}; \emph{Cordyceps militaris} & Iridium-coated stainless-steel subdermal & 1\,Hz (600 per s, averaged) & Resampled to 1\,Hz & SAX language encoding of spike trains; \emph{S.\ commune} most complex\\
\addlinespace
Olsson \& Hansson \cite{olsson}, 1995 & \emph{P.\ ostreatus}; \emph{Armillaria bulbosa} & Glass microelectrode & Not stated & Electrometer amplification & AP frequency increases with negative current; unlike animal neurons\\
\bottomrule
\end{tabularx}
\end{table}

\begin{instructions}
Retyped from the plan's image of this table. Consider adding a final row for this thesis: oyster (\emph{P.\ ostreatus},
three 16-channel recordings at 1\,Hz, about 1,500\,h with the held-out M4, plus its chocolate and snow oyster strains),
lion's mane (MJu26a, LM), Asian shiitake and reishi from the same lab; native rates of 1--10\,Hz kept; a detector sweep, seeds and classifiers
judged against blind verdicts.
\end{instructions}


A further trait of interest in the signal's behaviour is the observation that motifs of interest often occur in groupings, referred to as spike-trains. This behaviour has led to comparisons with the patterns seen in neuronal EEG recordings and cardiac ECG recordings \cite{slayman}, \cite{adamatzky2023}. A language analysis by Adamatzky \cite{adamatzky2022}, for example, treated these spike-trains as symbolic sequences and found that \emph{Schizophyllum commune} generates the most complex ``sentences'' compared to three other species when analysed at 1\,Hz, suggesting that there is meaningful structure in the timing and grouping of these events.

Note that there are many studies which examine the biological mechanisms of the measured potential differences, suggesting a propagation of calcium ions Ca\textsuperscript{+}, Ca\textsuperscript{2+} and potassium ions K\textsuperscript{+} through hyphae cell membranes both in general function and in response to stressors \cite{brand,itani,houdinet}. However, while this background is important for providing context about the potential origins of both the signal's shape and behaviours, this project is not concerned with the intracellular mechanisms that generate these signals, as it is deemed unnecessary for developing a comprehensive analysis pipeline and outside of the proposed scope.

A main drawback in this field is that it is relatively small, especially in comparison to the quantity of research in parallel fields on EEG and ECG analysis. This is most likely due to the fact that there is as of yet no clear financial incentive for investors to fund more extensive research into fungi signalling behaviour. As a result of this, there are many researchers who appear commonly across papers, most notably Andrew Adamatzky and Matteo Buffi. This implies that there is a certain risk in the current findings presented, as the biases within the work can be traced to only a handful of parties involved. A further consequence of this is that none of the analysis tools described in these studies have been released as public software, and preprocessing procedures are rarely described in sufficient detail to reproduce results independently.

\begin{instructions}
Qualify the last sentence with what this project found. Adamatzky's four-species recordings are public on Zenodo
\cite{adamatzky2021data}, and Dehshibi and Adamatzky released MATLAB code \cite{dehshibi2020code}. Porting that code
showed it differs from the printed paper in about ten places (Q33). Keep the point, that preprocessing and parameters
are rarely reproducible from the papers, but be precise about which tools are and are not released.
\end{instructions}

\section{The absence of a shared signal vocabulary}
\label{sec:litvocab}

The closest parallel domain to mycelium signal analysis is Electroencephalogram (EEG) and Electrocardiogram (ECG) analysis, in which a large body of methodological work has been possible precisely because a canonical event taxonomy was established early --- canonical sleep stages, R-peaks, and wave types are defined consistently enough across laboratories that automated classifiers can be trained on one dataset and evaluated on another \cite{murungi} \cite{mohammed}. The abundance of labelled datasets has allowed the machine-learning advances discussed in Section~\ref{sec:litnn}. Mycelium research lacks this foundation entirely, and the absence of a shared labelling standard is most likely a significant contributor to the stall in progress noted in Section~\ref{sec:litfungi}.

The consequence is concrete and is worth stating plainly, because it motivates the structure of this thesis. A comprehensive benchmark of modern time-series classification methods is provided by Bagnall et al.\ \cite{bagnall}, who compared 18 algorithms across a large suite of standard datasets and found that the best performing systems were ensemble and transformer-based classifiers that combined multiple feature representations. An important observation for the present project is that all benchmark datasets in that study were fully labelled prior to testing. A classifier cannot be trained or fairly evaluated without labels, and as such unsupervised discovery methods must precede any supervised approach. Where no ground truth exists, the ordering of work is forced: discovery first, then labelling, then classification --- and the labels themselves must be treated as a claim to be tested rather than as given.

\section{Prior research at UQ}
\label{sec:lituq}

Before beginning this thesis, I participated in the Underground Brains research project at The University of Queensland (see \url{https://eecs.uq.edu.au/project/fungi}) as part of a summer research program. The research involved the initial analysis of about 1,000 hours of bio-electric recordings on \emph{Pleurotus ostreatus}, more commonly known as oyster mushrooms. The recordings were made with 16 electrodes distributed across the fruiting bodies, sharing a common ground, and logged at 1\,Hz. The mushrooms were grown inside Faraday cages to limit outside interference.

The research was largely exploratory, and involved the development of a pipeline for effective data analysis in \textbf{MATLAB}. Using this pipeline, experimental analysis was performed which included ratio encoding and discovery, Symbolic Aggregate Approximation (SAX) encoding and discovery, Gramian Angular Field encoding, and wavelet and frequency analysis. The full repository developed can be found on Github \cite{summerrepo}. This experience involved developing a familiarity with the behaviour evident in the time-series recordings, early research into the field of time-series analysis and the initial implementations of several algorithms for analysis and motif extraction. While the current thesis builds on the knowledge developed here, particularly for the techniques in Gramian Angular Field Imaging as discussed in Section~\ref{sec:litnn}, the files, code and algorithms remain separated and unique, and can be found in a separate repository \cite{thesisrepo}.

The goal of UQ's Underground Brains project is to study the electromagnetic communication system of mycelium. As part of that, there are four main areas of research focus:

\begin{enumerate}
  \item Designing next generation sensors for environmental sensing
  \item Developing new recording and analysis techniques for the electromagnetic properties of mycelium
  \item Investigation of the communication and multi scale dynamics of the underground networks of mycelium
  \item Bridging art, science and engineering to help contribute to the public understanding of AI, STEM and fungi
\end{enumerate}

The motivation behind the project lies in the importance of fungi in the growth and stability of ecosystems. They are symbiotic with 90\% of plants, particularly in the diverse climates of Australia. However, fungi, in particular Australian species, have many distinct bio-chemical and electro-magnetic properties which are not well understood. Advances in how they communicate, grow and interact with their surrounding ecosystems has vital applications for Australian agriculture and ecosystem health.

The present thesis contributes primarily to the second of these areas, by developing the analytical tools and pipeline infrastructure needed to process and interpret the electromagnetic recordings described above.

\begin{instructions}
\textbf{Framing (5 Nov):} the summer project is separate prior work and must read that way. The line to hold: in the
summer, the SAX idea was explored; this year, the SAX algorithms (and everything else in Chapter~\ref{ch:early}) were
implemented anew in a separate code base. Do not describe any summer result as a thesis result.

\textbf{The 1,000 hours (5 Nov):} M2\_concat is oyster, so ``about 1,000 hours of \ldots\
\emph{Pleurotus ostreatus}'' above is right: M2\_aug 721\,h and M2\_concat 282.5\,h. The held-out M4 adds 505\,h of
oyster. Say which other recordings from the same lab joined the corpus during the year: lion's mane (MJu26a, LM),
chocolate oyster (CO) and snow oyster (SO), both strains of \emph{P.\ ostreatus}, Asian shiitake (SH) and reishi (F2A).
Chapter~\ref{ch:method} has the table.
\end{instructions}

\section{Gaps in knowledge}
\label{sec:gaps}

In today's information age, time-series analysis is a quickly growing field with many exciting concepts for extracting information from ever-more complex dynamical systems. However little to no research has explored the impact of multi-technique analysis, matrix profiling methods or machine-learning methods on the bio-electric signals of mycelium networks. These methods are fairly recent innovations in comparison to the scope of research on mycelium networks, and biological systems as a whole, and as such it is expected that their implementation into these fields has a period of lag. Existing examples of their implementations in similar fields examining non-stationary, multi-scalar time-series data has yielded exciting results \cite{shi} \cite{ponganis} \cite{liu}, suggesting that there is value to be found in these algorithms.

A second significant area in mycelium research which requires more attention is the idea of a standard of classifying and clustering similar signal families and similar motif structures. While similar motifs have been identified across papers and across species, no significant attempt has been made to generalise and group these due to their variability, and appearance across multiple time-scales. The ad hoc nature of existing labelling --- typically manual, species-specific, and non-transferable --- means that analysis results cannot be compared across datasets or used to train generalisable predictive models. Before these signals can be interpreted as separate events relating to significant biological functions, there must first be a system to extract and separate them automatically from a set of raw time-series data.

A third gap is the lack of any shared, publicly available computational pipeline for analysing these signals. The tools described in existing literature are typically bespoke, undocumented, and not released alongside papers \cite{dehshibi,buffi2025review}. This creates a reproducibility problem and prevents incremental progress from building on a shared foundation. Promising results are discovered and then lost, because the parameters and ordering that produced them were never captured as a reusable object.

\begin{instructions}
Consider adding a fourth gap from the project's own evidence: published detectors cannot be compared because they
disagree by orders of magnitude on the same recording, and no study reports a detector's precision against blind human
judgement. Chapter~\ref{ch:rq1} measures exactly this.

Also note the Dehshibi finding. The authors' released MATLAB differs from the printed paper in about ten places, and the
printed Algorithm 3 read literally yields inverted regions (\src{Working/Detection/analysis/dehshibi_authors.py},
QUESTIONS.md Q33). This is concrete evidence for the reproducibility gap.
\end{instructions}

\section{Research questions}
\label{sec:rqs}

The overall question this thesis addresses is unchanged from the thesis plan:

\begin{quote}
\textbf{Can the recurring structures present in mycelium bio-electric signals be automatically detected and characterised using modern time-series analysis methods?}
\end{quote}

The thesis plan (August 2026) approached this question through six sub-questions. As the instrument was built and the first results came in, those questions were refined and scoped to the three research questions below. Each is a question about what the instrument makes possible, and each is tested against blind human judgement and a matched null:

\begin{description}
  \item[RQ1 --- Detectors at scale.] Can multiple detection algorithms be compared at scale, across recordings, against blind human judgement and matched surrogates? And does agreement between independent detectors predict what a human accepts? (Chapter~\ref{ch:rq1})
  \item[RQ2 --- Seeds.] Can a human-checked motif from the library, used as a seed, find new instances that a human judging blind also accepts, more often than duration-matched random windows and far more often than the same search finds in a surrogate? (Chapter~\ref{ch:rq2})
  \item[RQ3 --- Classifiers.] Can labels gathered through the tool train a classifier that works on recordings it never saw, beating a label-shuffle null and the detectors' own vote? (Chapter~\ref{ch:rq3})
\end{description}

Two further chapters use the instrument to describe what it found. Chapter~\ref{ch:events} gives an overview of the events it found: their shape, size, scale and timing. Chapter~\ref{ch:roundb} analyses them through a set of shorter questions about the nature of the signals.

\begin{instructions}
Rewritten 5 Nov: the RQs are stated as they stand, with one sentence on their refinement from the August plan. The
plan-to-thesis mapping table and diagram moved to Chapter~\ref{ch:discussion}'s critical review
(Table~\ref{tab:rqmap}). Reword the two bridging sentences above in your own voice.
\end{instructions}

\section{Contributions}
\label{sec:contributions}

\begin{instructions}
New section, written by you. Suggested list:
\begin{enumerate}
  \item \tool{}, a composable, human-in-the-loop application for long unlabelled bio-electric recordings: 59 typed
        analysis blocks, 24 templates, a recipe hash on every result, surrogate nulls, a held-out lock, and an HPC route.
        Machine detections and human verdicts are stored apart.
  \item The first comparison of fungal spike and event detectors against blind, stratified human judgement at scale:
        1,960 verdicts across five recordings, with a detector-agreement ladder that holds on every recording.
  \item Seeded motif search shown to recover human-accepted events far above chance. Judged blind, 20/20 to 100/100
        of each seed's matches were accepted, against 1/12 to 9/40 of length-matched random windows (Fisher
        $p<0.0001$ for each seed).
  \item A classifier trained only on the tool's verdicts that transfers to unseen recordings (AUC 0.97 leave-one-out)
        and to the held-out recording (M4: AUC 0.98, balanced accuracy 0.93, scored once).
  \item A characterisation of accepted events: drops dominate; amplitude and duration ranges; scale collapse across
        three species; the finding that normalised shape still encodes resolution and species.
  \item A transferable method: recipes, nulls and a 342-page user manual (cited, not appended).
\end{enumerate}
\end{instructions}

\section{Scope and position on interpretation}
\label{sec:position}

This thesis measures recurrence and structure. It makes no claim about meaning, signalling, or communication.

This is a deliberate position rather than an omission. The framing of fungal electrical activity as ``communication'' or as a ``language'' is a live methodological criticism within the field \cite{buffi2025review}, and interpretive claims have in several cases run ahead of the evidence that recurring structure exists at all in a form that survives a statistical control. The order of operations adopted here is the opposite: confirm that patterns recur, at rates that exceed what a matched surrogate produces, and describe them in terms that another researcher could apply to their own recording --- before attempting to decode what they might correspond to biologically.

Three consequences follow, and they shape the whole document:
\begin{itemize}
  \item Every reported detection result is set against a null. For runs, that is the identical chain rerun on a phase-randomised or otherwise matched surrogate of the same signal (Section~\ref{sec:surrogates}). For human acceptance rates, it is a blind base rate.
  \item Cross-channel co-occurrence is not treated as evidence of anything until it has been separated from common-mode contamination through the shared ground, and from single propagating events (Section~\ref{sec:crosschannel}).
  \item The morphology vocabulary developed over the course of this project, the shape names in Appendix~\ref{app:vocab}, is treated as a hypothesis under test, not as ground truth.
\end{itemize}

Biological interpretation is out of scope. It is recorded in Section~\ref{sec:futurework} as future work, tied to the effectiveness of forecasting: a structure that can be used to predict what follows it is a considerably better candidate for interpretation than one that merely repeats.

\section{Overview of the thesis}

Chapter~\ref{ch:litreview} reviews the time-series methods used here and sets out the theory they depend on. Chapter~\ref{ch:early} describes the early experimental work of the year --- hand-labelled windows, image-encoded CNNs, matrix profiles, window matrices, symbolic encodings and a cross-species study of drops --- and what each taught. Chapter~\ref{ch:pipeline} describes \tool{}, the instrument those lessons led to: its design, how it was built and how it was evaluated. Chapter~\ref{ch:method} describes the corpus, the judging protocol, the controls and the held-out exam. Chapters~\ref{ch:rq1}, \ref{ch:rq2} and~\ref{ch:rq3} answer RQ1 to RQ3. Chapter~\ref{ch:events} gives an overview of the events the instrument found, and Chapter~\ref{ch:roundb} analyses them through a set of shorter questions about the signals. Chapter~\ref{ch:discussion} draws the results together, reviews the project critically against its plan and schedule, and sets out what should be done next.
